\documentclass[aps,prc,preprint,superscriptaddress,nofootinbib]{revtex4-2}

\usepackage{graphicx}
\usepackage{amsmath}
\usepackage{amssymb}
\usepackage{hyperref}
\usepackage{booktabs}
\usepackage{multirow}

\begin{document}

\title{Directed Flow of Light Nuclei in Au+Au Collisions at $\sqrt{s_{NN}} = 3.0$--$19.6$~GeV from STAR}

\author{Ding Chen}
\email{dchen087@ucr.edu}
\affiliation{University of California, Riverside}

% \author{Richard Seto}
% \affiliation{University of California, Riverside}

\date{\today}

\begin{abstract}
This analysis note describes the measurement of directed flow ($v_1$) for protons and light nuclei (deuterons, tritons, $^3$He, $^4$He) in Au+Au collisions at $\sqrt{s_{NN}} = 3.0$, 3.2, 4.5, and 7.7~GeV in fixed-target mode and 7.7 and 19.6~GeV in collider mode at STAR. The first-order event plane is reconstructed from the Event Plane Detector (EPD) using the three-subevent method with re-centering and Fourier shift corrections. Results are presented as $v_1/A$ vs.\ rapidity and $v_1/A$ vs.\ $p_T/A$ and compared with JAM2+coalescence model calculations. The atomic mass number ($A$) scaling predicted by the nucleon coalescence model is tested as a function of rapidity and collision energy. We observe approximate $A$-scaling at low energies near target rapidity, with systematic deviations emerging at mid-rapidity and higher energies.
\end{abstract}

\maketitle
\tableofcontents

%==============================================================================
\section{Introduction}
\label{sec:intro}
%==============================================================================

Directed flow ($v_1$), the first Fourier coefficient of the particle azimuthal distribution with respect to the reaction plane, is sensitive to the early-time dynamics of heavy-ion collisions. For light nuclei, directed flow provides a direct test of the production mechanism: the nucleon coalescence model predicts that $v_1$ scales linearly with atomic mass number $A$,
\begin{equation}
    v_1^{(A)} \approx A \cdot v_1^{(\text{nucleon})},
    \label{eq:ascaling}
\end{equation}
so that $v_1/A$ should be species-independent at the same rapidity and $p_T/A$.

Previous measurements from STAR at $\sqrt{s_{NN}} = 3.0$~GeV showed approximate $A$-scaling for $v_1$ but clear breaking for $v_2$~\cite{STAR:2021yiu}. The HADES Collaboration reported deviations from simple $A$-scaling in the few-GeV regime~\cite{HADES:2020lob}. Measurements at $\sqrt{s_{NN}} = 7.7$--27~GeV (BES-II, collider mode) show that light nuclei $v_1$ slopes are positive and follow a smooth energy dependence, with no clear $A$-scaling observed~\cite{Wang:2026}.

This analysis extends the $v_1$ measurements of light nuclei across the STAR FXT energy range ($\sqrt{s_{NN}} = 3.0$--7.7~GeV) and collider energies (7.7, 19.6~GeV), providing a systematic study of $A$-scaling from high baryon density (target rapidity at low energies) to lower baryon density (mid-rapidity at higher energies). Comparisons with JAM2 transport model calculations incorporating nucleon coalescence provide theoretical context.

%==============================================================================
\section{Dataset and Event Selection}
\label{sec:dataset}
%==============================================================================

\subsection{Experimental Setup}
\label{subsec:setup}

In fixed-target (FXT) mode, a 1~mm thick gold foil is placed at $z = 200$~cm inside the STAR TPC. A single Au beam strikes the stationary target, producing center-of-mass energies $\sqrt{s_{NN}} = 3.0$, 3.2, 4.5, and 7.7~GeV. The collision products travel predominantly forward (toward negative $z$ in STAR coordinates) and are detected by the Time-Projection Chamber (TPC), Time-of-Flight (TOF) detector, and the east-side Event Plane Detector (EPD).

In collider (COL) mode, two counter-circulating Au beams collide near the center of the detector ($z \approx 0$). We analyze collider datasets at $\sqrt{s_{NN}} = 7.7$ and 19.6~GeV from the BES-II program (2019--2021).

Table~\ref{tab:rapidity_params} lists the kinematic parameters for each collision energy.

\begin{table}[h]
\centering
\caption{Kinematic parameters for each collision energy.}
\label{tab:rapidity_params}
\begin{tabular}{ccccc}
\toprule
$\sqrt{s_{NN}}$ (GeV) & Mode & $y_{\text{mid}}$ & $y_{\text{beam,CM}}$ & Year \\
\midrule
3.0  & FXT & $-1.05$  & $\pm 1.05$  & 2018 \\
3.2  & FXT & $-1.14$  & $\pm 1.14$  & 2019 \\
4.5  & FXT & $-1.52$  & $\pm 1.52$  & 2020 \\
7.7  & FXT & $-2.10$  & $\pm 2.10$  & 2020 \\
7.7  & COL & 0        & $\pm 2.09$  & 2021 \\
19.6 & COL & 0        & $\pm 3.04$  & 2019 \\
\bottomrule
\end{tabular}
\end{table}

\subsection{Sign Convention}
\label{subsec:sign_convention}

In FXT mode, the beam travels toward negative $z$ (eastward), so $y_{\text{beam,lab}} < 0$. Center-of-mass rapidity is $y_{CM} = y_{\text{lab}} - y_{\text{mid}}$. In the CM frame, positive $y_{CM}$ corresponds to the target rapidity direction and negative $y_{CM}$ to the beam rapidity direction. We maintain the same sign conventions as established STAR collider-mode analyses to facilitate direct comparison.

\subsection{Event Selection}
\label{subsec:event_selection}

Only minimum-bias triggered events are used. Events must have a well-reconstructed primary vertex. Table~\ref{tab:dataset_summary} summarizes the trigger IDs, vertex cuts, and event statistics for each dataset.

\begin{table}[h]
\centering
\caption{Dataset summary: trigger IDs, vertex cuts, and event statistics.}
\label{tab:dataset_summary}
\begin{tabular}{cccccc}
\toprule
$\sqrt{s_{NN}}$ (GeV) & Mode & Trigger ID & $V_z$ (cm) & $V_r$ (cm) & Good Events \\
\midrule
3.0  & FXT & 620052 & $[198, 202]$  & $< 1.5$ & $256 \times 10^{6}$ \\
3.2  & FXT & 680001 & $[198, 202]$  & $< 1.5$ & $159 \times 10^{6}$ \\
4.5  & FXT & 780001 & $[198, 202]$  & $< 1.5$ & [TBD] \\
7.7  & FXT & 780020 & $[198, 202]$  & $< 1.5$ & [TBD] \\
7.7  & COL & 810010--810040 & $|V_z| < 70$   & $< 2.0$ & [TBD] \\
19.6 & COL & 640001--640051 & $|V_z| < 70$   & $< 2.0$ & $443 \times 10^{6}$ \\
\bottomrule
\end{tabular}
\end{table}

Runs with abnormal detector performance are excluded based on the official STAR bad run lists for each dataset.

\subsection{Centrality Definition}
\label{subsec:centrality}

Centrality is determined from the reference multiplicity (number of charged tracks within a defined pseudorapidity range). The official STAR centrality definitions are used for each energy. Due to low trigger efficiency for peripheral collisions in FXT mode, the analysis uses centrality 0--60\%. The primary centrality range for flow results is 10--40\%.

%==============================================================================
\section{Track Quality Cuts}
\label{sec:track_cuts}
%==============================================================================

Reconstructed tracks must satisfy quality criteria to ensure reliable momentum and energy-loss measurements. Table~\ref{tab:track_cuts} lists the selection criteria applied uniformly across all energies.

\begin{table}[h]
\centering
\caption{TPC track quality cuts.}
\label{tab:track_cuts}
\begin{tabular}{lc}
\toprule
Variable & Requirement \\
\midrule
nHitsFit       & $\geq 15$ \\
nHitsDedx      & $> 5$ \\
nHitsFit/nHitsPoss & $> 0.52$ \\
DCA            & $< 3.0$~cm \\
\bottomrule
\end{tabular}
\end{table}

The nHitsFit cut ensures sufficient tracking points for reliable momentum reconstruction. The nHitsFit/nHitsPoss ratio cut rejects split tracks from particle decays. The distance of closest approach (DCA) cut selects primary tracks originating from the collision vertex.

%==============================================================================
\section{Particle Identification}
\label{sec:pid}
%==============================================================================

\subsection{TPC $dE/dx$ and $n\sigma$}
\label{subsec:nsigma}

The TPC measures specific ionization energy loss ($dE/dx$) for each track. The deviation from the expected Bichsel parametrization~\cite{Bichsel:2006cs} for particle species $X$ is quantified as:
\begin{equation}
    n\sigma_X = \frac{1}{R} \ln\left(\frac{(dE/dx)_{\text{meas}}}{(dE/dx)_{X,\text{exp}}}\right),
    \label{eq:nsigma}
\end{equation}
where $R$ is the $dE/dx$ resolution. Protons are identified with $|n\sigma_p| < 2$. For light nuclei with $Z \geq 2$ ($^3$He, $^4$He), the $dE/dx$ bands are well separated due to the $Z^2/\beta^2$ scaling of energy loss.

\subsection{TOF Mass-Squared}
\label{subsec:m2}

The Time-of-Flight detector measures the velocity $\beta$ of each particle, allowing computation of the squared mass:
\begin{equation}
    m^2 = p^2\left(\frac{1}{\beta^2} - 1\right).
    \label{eq:m2}
\end{equation}

\subsection{PID Cuts by Species}
\label{subsec:pid_cuts}

Table~\ref{tab:pid_cuts} summarizes the PID criteria for each particle species.

\begin{table}[h]
\centering
\caption{PID selection criteria for each particle species.}
\label{tab:pid_cuts}
\begin{tabular}{lccc}
\toprule
Species & $n\sigma$ cut & $m^2$ (GeV$^2/c^4$) & TOF req. \\
\midrule
$p$       & $|n\sigma_p| < 2.0$ & ---                & No  \\
$d$       & $|n\sigma_d| < 2.0$ & $2.8 < m^2 < 4.5$  & Yes \\
$t$       & $|n\sigma_t| < 2.0$ & $6.0 < m^2 < 10.0$ & Yes \\
$^3$He    & $|n\sigma_{^3\text{He}}| < 3.0$ & ---     & No  \\
$^4$He    & $|n\sigma_{^4\text{He}}| < 3.0$ & ---     & No  \\
\bottomrule
\end{tabular}
\end{table}

%==============================================================================
\section{Event Plane Reconstruction}
\label{sec:event_plane}
%==============================================================================

\subsection{Q-vector and Event Plane Angle}
\label{subsec:qvector}

The first-order event plane angle $\Psi_1$ is reconstructed from the flow vector:
\begin{equation}
    Q_{1,x} = \sum_i w_i \cos\phi_i, \quad Q_{1,y} = \sum_i w_i \sin\phi_i,
\end{equation}
\begin{equation}
    \Psi_1 = \arctan\left(\frac{Q_{1,y}}{Q_{1,x}}\right),
\end{equation}
where $\phi_i$ is the azimuthal angle and $w_i$ is the weight of particle/hit $i$. For TPC tracks, $w_i = p_T$. For EPD tiles, $w_i = \text{TnMIP}$ (truncated nMIP, threshold 0.3, maximum 2.0).

\subsection{Sub-Event Configuration}
\label{subsec:subevents}

In FXT mode, the forward and backward rapidity regions have different multiplicities and flow magnitudes. We define four subevents:
\begin{itemize}
    \item \textbf{EPD A}: Inner rings (rows 1--8), $-5.6 < \eta < -3.5$
    \item \textbf{EPD B}: Outer rings (rows 9--16), $-3.5 < \eta < -2.4$
    \item \textbf{TPC A}: $\eta < \eta_{\text{mid}} - 0.05$
    \item \textbf{TPC B}: $\eta > \eta_{\text{mid}} + 0.05$
\end{itemize}
The main event plane is reconstructed from EPD A, with EPD B and TPC B as reference subevents for resolution.

In COL mode, the EPD east and west halves and TPC halves divided at $\eta = 0$ serve as subevents.

\subsection{Re-centering and Fourier Shifting}
\label{subsec:corrections}

Detector acceptance non-uniformities are corrected in two steps:

\textbf{Re-centering}: The Q-vector components are shifted by their event-averaged means:
\begin{equation}
    Q'_{x} = Q_x - \langle Q_x \rangle, \quad Q'_{y} = Q_y - \langle Q_y \rangle.
\end{equation}

\textbf{Fourier shifting}: Residual modulations are flattened using:
\begin{equation}
    \Delta\Psi = \sum_{k=1}^{k_{\max}} \frac{2}{k} \left[-\langle\sin(k\Psi)\rangle \cos(k\Psi) + \langle\cos(k\Psi)\rangle \sin(k\Psi)\right],
\end{equation}
with $k_{\max} = 10$. After correction, the event plane distributions are uniform in $[-\pi, \pi]$.

\subsection{Event Plane Resolution}
\label{subsec:resolution}

In FXT mode, the three-subevent method is used since the subevents have unequal resolution:
\begin{equation}
    R_{11}^A = \sqrt{\frac{\langle\cos(\Psi_1^A - \Psi_1^B)\rangle \langle\cos(\Psi_1^A - \Psi_1^C)\rangle}{\langle\cos(\Psi_1^B - \Psi_1^C)\rangle}},
    \label{eq:resolution_3sub}
\end{equation}
where $A$, $B$, $C$ denote the three subevents (EPD A, EPD B, TPC B).

In COL mode, the two-subevent method is used. The subevent resolution $R_{1,\text{sub}} = \sqrt{\langle\cos(\Psi_1^A - \Psi_1^B)\rangle}$ is converted to the full resolution by numerically inverting the Voloshin function~\cite{Poskanzer:1998yz}:
\begin{equation}
    R_k(\chi) = \frac{\sqrt{\pi}}{2\sqrt{2}}\,\chi\,e^{-\chi^2/4}\left[I_{(k-1)/2}\!\left(\frac{\chi^2}{4}\right) + I_{(k+1)/2}\!\left(\frac{\chi^2}{4}\right)\right],
\end{equation}
to find $\chi_{\text{sub}}$, then computing $R_{1,\text{full}} = R_1(\sqrt{2}\,\chi_{\text{sub}})$.

%==============================================================================
\section{Flow Measurement}
\label{sec:flow}
%==============================================================================

The directed flow is measured as:
\begin{equation}
    v_1 = \frac{\langle\cos(\phi - \Psi_1)\rangle}{R_{11}},
    \label{eq:v1}
\end{equation}
where $\phi$ is the particle azimuthal angle and $R_{11}$ is the event plane resolution. Each particle is weighted by $1/\epsilon$, where $\epsilon = \epsilon_{\text{TPC}} \cdot \epsilon_{\text{TOF}}$ is the combined efficiency correction.

Results are presented as $v_1/A$ vs.\ $y_{CM}$ and $v_1/A$ vs.\ $p_T/A$ in rapidity windows defined in Table~\ref{tab:rapidity_windows}.

\begin{table}[h]
\centering
\caption{Rapidity windows for $p_T/A$-dependent analysis (in $y_{CM}$). Values are energy-dependent and correspond to approximately equal fractions of $y_{\text{beam,CM}}$.}
\label{tab:rapidity_windows}
\begin{tabular}{ccccc}
\toprule
$\sqrt{s_{NN}}$ (GeV) & yMid & yFor & yExt & yEnd \\
\midrule
3.0  & $[0.1, 0.25]$ & $[0.25, 0.5]$ & $[0.5, 0.75]$  & $[0.75, 1.0]$  \\
3.2  & $[0.1, 0.25]$ & $[0.25, 0.5]$ & $[0.5, 0.75]$  & $[0.75, 1.1]$  \\
4.5  & $[0.15, 0.38]$ & $[0.38, 0.76]$ & $[0.76, 1.14]$ & $[1.14, 1.52]$ \\
7.7  & $[0.1, 0.53]$ & $[0.53, 1.05]$ & $[1.05, 1.58]$ & $[1.58, 2.10]$ \\
\bottomrule
\end{tabular}
\end{table}

%==============================================================================
\section{Efficiency Corrections}
\label{sec:efficiency}
%==============================================================================

\subsection{TPC Tracking Efficiency}

TPC reconstruction efficiencies are determined from embedding simulations. Monte Carlo tracks are embedded into real events and processed through the full STAR reconstruction chain. The efficiency is:
\begin{equation}
    \epsilon_{\text{TPC}}(y_{CM}, p_T) = \frac{N_{\text{matched}}(y_{CM}, p_T)}{N_{\text{generated}}(y_{CM}, p_T)}.
\end{equation}
Separate efficiency maps are produced for each particle species at each energy.

\subsection{TOF Matching Efficiency}

TOF matching efficiency is extracted from data:
\begin{equation}
    \epsilon_{\text{TOF}}(y_{CM}, p_T) = \frac{N_{\text{TPC+TOF}}(y_{CM}, p_T)}{N_{\text{TPC}}(y_{CM}, p_T)}.
\end{equation}

Each particle is weighted by $1/(\epsilon_{\text{TPC}} \cdot \epsilon_{\text{TOF}})$ when filling flow histograms.

%==============================================================================
\section{Systematic Uncertainties}
\label{sec:systematics}
%==============================================================================

Systematic uncertainties are estimated by varying analysis cuts. The following sources are considered:

\begin{itemize}
    \item DCA cut: varied from 1.0 to 3.0~cm (nominal: 3.0~cm)
    \item nHitsFit: varied from 15 to 20 (nominal: 15)
    \item nHitsFit/nHitsPoss: varied from 0.52 to 0.60
    \item $n\sigma$ width: $\pm 30\%$
    \item $m^2$ window: $\pm 30\%$
    \item Vertex cuts: $V_z$ and $V_r$ ranges varied
\end{itemize}

The Barlow check~\cite{Barlow:2002yb} identifies significant variations:
\begin{equation}
    \frac{|\Delta v_1|}{\sqrt{|\sigma_{\text{var}}^2 - \sigma_{\text{nom}}^2|}} > 1.
    \label{eq:barlow}
\end{equation}
The systematic uncertainty in each bin is the standard deviation of all significant variations.

%==============================================================================
\section{Model Comparison: JAM2 + Coalescence}
\label{sec:model}
%==============================================================================

Data are compared with JAM2~\cite{Nara:2019qfd}, the C++ implementation of the JAM hadronic transport model, using:
\begin{itemize}
    \item RQMD.RMF mean field (momentum-dependent, MD2 parametrization)
    \item Built-in nucleon coalescence: $R_0 = 3.0$~fm, $P_0 = 0.25$~GeV/$c$
    \item $\sim 10^6$ events per energy
\end{itemize}

Centrality is defined by impact parameter cuts. Flow is computed using the known reaction plane.

Note: the JAM1 afterburner coalescence uses different parameters ($R_0 = 3.575$~fm, $P_0 = 0.285$~GeV/$c$), resulting in a larger coalescence phase-space volume. A comparison of both parameter sets is planned.

Code: JAM1 (\url{https://github.com/ynara305/jam1}), JAM2 (\url{https://gitlab.com/transportmodel/jam2}).

%==============================================================================
\section{Results}
\label{sec:results}
%==============================================================================

Results are for 10--40\% centrality unless noted.

\subsection{$v_1/A$ vs.\ Rapidity}
\label{subsec:v1_vs_y}

% TODO: Insert figures for each energy

At $\sqrt{s_{NN}} = 3.0$ and 3.2~GeV, approximate $A$-scaling is observed: $v_1/A$ for deuterons overlaps with protons across the measured rapidity range. Near target rapidity ($y_{CM}/y_{\text{beam}} \approx 0.9$), both data and JAM2 show consistent $A$-scaling.

At $\sqrt{s_{NN}} = 4.5$~GeV, deviations from $A$-scaling begin to appear near mid-rapidity, while approximate scaling persists at forward rapidity.

At $\sqrt{s_{NN}} = 7.7$~GeV (FXT), proton and deuteron $v_1/A$ show clear separation at mid-rapidity, indicating breaking of simple $A$-scaling.

At $\sqrt{s_{NN}} = 19.6$~GeV (COL), both proton and deuteron $v_1$ slopes approach zero near mid-rapidity.

\subsection{$v_1/A$ vs.\ $p_T/A$}
\label{subsec:v1_vs_pt}

% TODO: Insert figures for each energy

The transverse momentum dependence of $v_1/A$ is examined in four rapidity windows (yMid, yFor, yExt, yEnd) at each energy. For species where $A$-scaling holds, the $v_1/A$ vs.\ $p_T/A$ curves should overlap across species.

\subsection{Energy Dependence of $A$-Scaling}
\label{subsec:energy_dependence}

% TODO: Insert summary figure

The degree of $A$-scaling is quantified by comparing the $v_1$ slope $dv_1/dy|_{y=0}$ normalized by $A$. Deviations grow with collision energy and decrease toward target rapidity.

%==============================================================================
\section{Summary}
\label{sec:summary}
%==============================================================================

We present measurements of directed flow ($v_1$) for protons, deuterons, tritons, $^3$He, and $^4$He in Au+Au collisions at $\sqrt{s_{NN}} = 3.0$--19.6~GeV. Key findings:
\begin{itemize}
    \item Approximate $A$-scaling of $v_1$ at $\sqrt{s_{NN}} = 3.0$ and 3.2~GeV across measured rapidity.
    \item Deviations from $A$-scaling at $\sqrt{s_{NN}} \geq 4.5$~GeV, particularly at mid-rapidity.
    \item $A$-scaling better preserved at forward (target) rapidity at all energies.
    \item JAM2+coalescence reproduces qualitative trends; quantitative agreement depends on coalescence parameters.
    \item $p_T/A$ dependence provides additional discriminating power for the coalescence mechanism.
\end{itemize}

\begin{thebibliography}{99}

\bibitem{STAR:2021yiu}
M.~S.~Abdallah \textit{et al.} (STAR Collaboration),
Phys.\ Rev.\ C \textbf{104}, 014908 (2021).

\bibitem{HADES:2020lob}
J.~Adamczewski-Musch \textit{et al.} (HADES Collaboration),
Phys.\ Rev.\ Lett.\ \textbf{125}, 262301 (2020).

\bibitem{Wang:2026}
M.~M.~Wang (STAR Collaboration),
Analysis note, January 2026.

\bibitem{Bichsel:2006cs}
H.~Bichsel,
Nucl.\ Instrum.\ Meth.\ A \textbf{562}, 154 (2006).

\bibitem{Poskanzer:1998yz}
A.~M.~Poskanzer and S.~A.~Voloshin,
Phys.\ Rev.\ C \textbf{58}, 1671 (1998).

\bibitem{Barlow:2002yb}
R.~Barlow,
hep-ex/0207026 (2002).

\bibitem{Nara:2019qfd}
Y.~Nara, T.~Maruyama, and H.~Stoecker,
Phys.\ Rev.\ C \textbf{102}, 024913 (2020).

\bibitem{Butler:1963}
S.~T.~Butler and C.~A.~Pearson,
Phys.\ Rev.\ \textbf{129}, 836 (1963).

\end{thebibliography}

\end{document}
